% Week 5 summary of thesis/risk_models/01_model_features_v2.tex. Numbers and citations are carried over
% from that file unchanged; nothing new was read for this note.
\documentclass[11pt,a4paper]{article}
\usepackage[margin=1in]{geometry}
\usepackage{amsmath}
\usepackage{booktabs}
\usepackage{array}
\usepackage{microtype}
\usepackage[colorlinks=true,linkcolor=blue,citecolor=blue,urlcolor=blue]{hyperref}

\title{Week 5: the plan for predicting new plaque from tree geometry}
\author{}
\date{}

\begin{document}
\maketitle

\section*{Progress}

\textbf{Question fixed.} The overall research question this thesis sets out to answer is: Does geometry of a coronary tree that is still healthy at the first scan predict where new plaque forms at the second scan? Former research found that low wall shear stress is a risk factor for plaque formation, and geometry has been proven to predict wall shear. Therefore is is a plausable theory to test that geometry can predict plaque risk. For the most part research has been conducted on trees with plaque present. 

The question is now phrased in terms of geometry alone, because the thesis will not simulate flow through every tree of a population cohort.
The data are the Herlev--{\O}sterbro (CGPS) cohort,
the outcome is new plaque in an artery that had none at baseline, and I have chosen features going into my risk predictions model has some geometric tie to wall shear. All features are computed from the centerline and the vessel radius alone.


\textbf{Features locked.} Three bifurcations are measured per tree: the left main (LM) into the left
anterior descending (LAD) and left circumflex (LCx), the LAD giving off its first diagonal (D1), and
the LCx giving off its first obtuse marginal (OM1). Trees with a ramus intermedius are flagged and
every analysis is repeated without them. A missing D1 or OM1 is coded as absent, not dropped.
Radii are the median Euclidean distance transform over a 3\,mm window that starts one junction
radius past the branch point.

\begin{center}
\small
\begin{tabular}{@{}p{3.3cm}p{5.6cm}p{5.6cm}@{}}
\toprule
Feature & Definition & Expected direction \\
\midrule
Angle $\theta$ & angle between the two daughters' fitted directions, 3\,mm window &  (wider angle, more risk) \\
Daughter ratio $\rho$ & $\min(r_1,r_2)/\max(r_1,r_2)$; used at D1, OM1 and for the LM outcome &  (lopsided split, more risk) \\
LM log ratio $\lambda$ & $\ln(r_\mathrm{LAD}/r_\mathrm{LCx})$; used for LAD and LCx outcomes & for LAD, for LCx \\
Deviation $\Delta_{HK}$ & $\frac{1}{r_0}\big(\frac{r_1^{7/3}+r_2^{7/3}}{2}\big)^{3/7}-2^{-3/7}$, the departure from the Huo--Kassab law & sign open, tested two-sided \\
Mean curvature $\bar\kappa_a$ & $\frac{1}{L}\int_0^L|\kappa(s)|\,\mathrm{d}s$ from the vessel origin & \\
Curvature profile $\kappa(s)$ & the curve itself, keeping where along the artery the bend sits & tested against the mean \\
\bottomrule
\end{tabular}
\end{center}

Ratios of radii were chosen because they cancel any multiplicative scale error, as in Finet et al.
Segmentation blur is additive and does not cancel, so its effect on each feature is measured on the
ImageCAS-X trees (one voxel of dilation and erosion, junction offset 0.5, 1 and 2 $r_J$, window 2, 3 and
5\,mm) before any outcome is opened. $\lambda$ and $\rho$ carry the same information at the LM, so
they never enter the same model.

\textbf{Metrics dropped.} The tortuosity index (explained no low-shear area in Kashyap et al.,
$R^2\approx 0$), root-mean-square curvature (same bending as the mean, with a noisier derivative), the
junction exponent (no closed form and biased by radius noise), the Finet ratio (an exact function of
$\rho$ and $\Delta_{HK}$, so its weight could not be interpreted), torsion (third derivative of a noisy
centerline; the median LAD is planar to within 2.1\% of its length yet the code reports a full twist every
12\,mm) and right coronary crux features (no supporting study, and they mostly re-encode dominance).
The finding on torsion comes from \texttt{scripts/check\_vessel\_planarity.py} on the 160 test cases.

\textbf{Four models.} Each adds one thing to the one before.
\begin{itemize}
  \item \textbf{Model 0:} logistic regression on age, sex and dominance, plus classical risk factors
    and the calcium score where recorded. The baseline that geometry has to beat.
  \item \textbf{Model 1:} adds the geometric features, with curvature as $\bar\kappa_a$.
  \item \textbf{Model 2:} replaces $\bar\kappa_a$ by $K$ (3 to 5) functional principal component scores of
    $\kappa(s)$, fitted on ImageCAS-X. The likelihood-ratio test against Model 1 is the primary test of
    whether the position of a bend adds anything to the mean.
  \item \textbf{Model 3:} a small one-dimensional convolutional network on the profile with late fusion of
    the scalar features, against a control network given $\bar\kappa_a$. Exploratory only, since it has far
    more parameters than the events can support.
\end{itemize}s
Age, sex and dominance enter every model, because what is normal depends on them (a left-dominant tree
has a larger LCx, so a lower normal $\lambda$).

\textbf{One model set per artery domain.} LAD, LCx, RCA and LM are each predicted on their own,
because shear acts locally and one score per patient would average a harmful geometry in one artery
with a harmless one in another. The RCA has no bifurcation feature, so it isolates curvature as a mean
against a profile. The LM is too short for a profile, so it gets Models 0 and 1 with $\rho$ in place of
$\lambda$. Models are fitted on participants whose whole tree was plaque free at baseline, because an
artery enlarges outward as plaque accumulates. Model 1 has 12 weights for LAD and LCx, 5 for RCA and 7
for LM.

\textbf{Statistics fixed before any outcome is opened.}
\begin{itemize}
  \item \emph{Sample size:} the Riley criteria, per domain and separately for the LM. Where events are
    too few, the side-branch block is penalised harder in all models. No events-per-weight rule is used.
  \item \emph{Fitting:} ridge penalty chosen by nested cross-validation.
  \item \emph{Validation:} the bootstrap resamples whole patients, since a patient contributes an outcome
    to every domain. Reported are the optimism-corrected AUC, calibration slope and intercept, and Brier
    score (TRIPOD+AI). Internal validation only, and this is stated with every result.
  \item \emph{Scan interval:} unknown onset time, so the interval enters as a covariate or the outcome is
    modelled in discrete time. Baseline features of participants who died, had an infarction or did not
    return are compared with those rescanned.
  \item \emph{Correlated features:} harmless for prediction, misleading for interpretation. Correlations
    and (generalised) VIFs are reported on ImageCAS-X first and on Herlev--{\O}sterbro after, with no VIF
    cutoff. Weights are read through ridge or one descriptor per family ($|r|>0.7$ flags a family), and
    family importance by likelihood-ratio test with a permutation reference under ridge. The side-branch
    features enter as $I\mathbf{x}$ with $\mathbf{x}$ centred, so $I$ and $I\mathbf{x}$ are uncorrelated.
  \item \emph{Model 3 importance:} refitting without a feature family, not shuffling (Hooker et al.).
\end{itemize}

\textbf{Locked on ImageCAS-X first.} The reliability floor for $\Delta_{HK}$, the length $L$ and the rule
for arteries shorter than $L$, $K$ for the functional PCA, and the correlation matrices.

\subsection*{Open points}

\begin{itemize}
  \item The size of the rescanned CGPS subset and the interval between scans are still to be
    confirmed once data access is granted. The sample-size calculation depends on both.
  \item Model 3 cannot be justified by the events, so a failure to beat Model 2 is not evidence against
    the profile.
  \item The Huo--Kassab deviation has no coronary study to lean on, so its direction is treated as open.
\end{itemize}

\end{document}
